\documentclass{article}
\usepackage[T1]{fontenc}
\usepackage{lmodern}
\usepackage{graphicx}
\usepackage{amsmath,amssymb}
\usepackage[a4paper,margin=2cm]{geometry}
\usepackage[absolute,overlay]{textpos}
\setlength{\TPHorizModule}{1mm}
\setlength{\TPVertModule}{1mm}
\usepackage[indonesian]{babel}
\usepackage{hyperref}
\setlength{\emergencystretch}{2em}
% unit: Halaman 4

\begin{document}

% === Halaman 4 (llm) ===

\section*{Kurva Eliptik}

\begin{itemize}
    \item Kurva eliptik adalah kurva dengan bentuk umum persamaan:
    \[ y^2 = x^3 + ax + b \]
    dengan syarat $4a^3 + 27b^2 \neq 0$
    \item Tiap nilai a dan b yang berbeda memberikan kurva eliptik yang berbeda pula.
\end{itemize}

\end{document}
